\documentclass[10pt,a4paper]{article}
\usepackage[margin=0.5in]{geometry}
\usepackage{enumitem}
\usepackage{titlesec}
\usepackage{hyperref}
\usepackage{xcolor}
\usepackage[T1]{fontenc}
\usepackage{lmodern}
\usepackage{tabularx}

\definecolor{linkblue}{RGB}{10,60,150}
\hypersetup{colorlinks=true, urlcolor=linkblue, linkcolor=linkblue}

\pagestyle{empty}
\setlength{\parindent}{0pt}
\setlength{\parskip}{0pt}

\titleformat{\section}{\large\bfseries\scshape}{}{0em}{}[\titlerule]
\titlespacing*{\section}{0pt}{5pt}{3pt}

\newcommand{\entry}[4]{%
  \textbf{#1} \hfill #2 \\
  \textit{#3} \hfill \textit{#4} \\[1pt]
}

\newcommand{\entrysingle}[2]{%
  \textbf{#1} \hfill \textit{#2} \\[2pt]
}

\begin{document}

% ===== HEADER =====
\begin{center}
  {\LARGE \textbf{Brhanu Fentaw Znabu}} \\[4pt]
  \href{mailto:bfentaw2@huskers.unl.edu}{bfentaw2@huskers.unl.edu} $\cdot$
  \href{mailto:brhanufenbme@gmail.com}{brhanufenbme@gmail.com} $\cdot$
  531-254-0990 $\cdot$ Lincoln, NE \\[2pt]
  \href{https://brhanufen.github.io}{brhanufen.github.io} $\cdot$
  \href{https://linkedin.com/in/brhanu-fentaw-znabu-23115619a}{LinkedIn} $\cdot$
  \href{https://github.com/brhanufen}{GitHub}
\end{center}

% ===== RESEARCH INTERESTS =====
\section{Research Interests}
My research develops deep learning foundation and generative models for large-scale biological sequence analysis across genomes and proteins, with a focus on domain-adaptive pre-training, generative design, and rigorous leakage-aware evaluation. I build systems end to end, from data curation through HPC training to open-source release, and apply them to predicting evolutionary trajectories and protein function.

% ===== EDUCATION =====
\section{Education}

\entry{Ph.D. Biomedical Engineering}{Aug 2024 -- Expected May 2029}
{University of Nebraska--Lincoln}{Lincoln, NE}
\begin{itemize}[leftmargin=1.2em, topsep=0pt, itemsep=1pt, parsep=0pt]
  \item Advisors: Nicole R. Sexton, PhD (Nebraska Center for Virology) $|$ Qiuming Yao, PhD (School of Computing)
\end{itemize}
\vspace{1pt}

\entry{M.Sc. Biomedical Science and Engineering}{2021 -- 2024}
{Gwangju Institute of Science and Technology (GIST)}{South Korea}
\begin{itemize}[leftmargin=1.2em, topsep=0pt, itemsep=1pt, parsep=0pt]
  \item GPA: 4.0/4.5 $|$ Advisor: Euiheon Chung, PhD
  \item Thesis: Deep behavioral phenotyping of diabetic neuropathy mouse model using open-field test
\end{itemize}
\vspace{1pt}

\entry{B.Sc. Biomedical Engineering}{2014 -- 2018}
{Jimma University}{Ethiopia}
\begin{itemize}[leftmargin=1.2em, topsep=0pt, itemsep=1pt, parsep=0pt]
  \item GPA: 3.64/4.0 $|$ Graduated with Distinction (Ranked 4th of 180)
\end{itemize}

% ===== RESEARCH EXPERIENCE =====
\section{Research Experience}

\entry{PhD Researcher}{Aug 2024 -- Present}
{Sexton Lab, University of Nebraska--Lincoln}{Lincoln, NE}
\begin{itemize}[leftmargin=1.2em, topsep=0pt, itemsep=1pt, parsep=0pt]
  \item Built ArboFM, a domain-adapted transformer model via continued masked language model pre-training of DNABERT-2 (117M params) on 120K+ sequence windows across 362 species; achieved AP = 0.978 under species-grouped cross-validation preventing data leakage.
  \item Designed a leakage-aware evaluation framework quantifying 16-pp performance inflation from phylogenetic data leakage in cross-validation; applied gradient-based attribution and 6-mer enrichment analysis for model interpretability.
  \item Built a cross-domain transfer learning matrix across four biological families (6,052 sequences, 124K+ windows); demonstrated that learned representations are domain-specific (Jaccard $\approx$ 0), informing transfer learning strategy design.
  \item Performed retrospective temporal validation demonstrating the model detects distribution shifts in held-out future data, including within-class lineage discrimination (CLES = 0.898).
  \item Developed ViraPredict, a multi-modal forecasting framework integrating fitness landscape modeling, ESM-2 (650M params) zero-shot mutation scoring (F1 = 0.936), and LoRA-fine-tuned DNABERT-2 with temporal embeddings on 29,493 curated sequences across 130 countries.
\end{itemize}
\vspace{1pt}

\entry{Research Assistant}{Mar 2021 -- Jun 2024}
{Neurophotonics Lab, GIST}{South Korea}
\begin{itemize}[leftmargin=1.2em, topsep=0pt, itemsep=1pt, parsep=0pt]
  \item Developed unsupervised autoregressive and Hidden Markov models for temporal sequence modeling and state transition detection in behavioral time-series data.
  \item Benchmarked deep learning-based 3D pose estimation (DeepLabCut) against traditional 2D analysis using logistic regression on multi-dimensional behavioral features.
\end{itemize}

% ===== OPEN-SOURCE PROJECTS =====
\section{Open-Source Projects}

\begin{itemize}[leftmargin=1.2em, topsep=2pt, itemsep=2pt, parsep=0pt]
  \item \textbf{TrustBinder} \href{https://github.com/brhanufen/trustbinder}{(github.com/brhanufen/trustbinder)}: Dual-oracle audit for de novo binder design, testing whether the score used to select generated binders survives an independent structure oracle. Boltz-2 and AlphaFold2 are uncorrelated on PD-L1 (r = 0.08, 100 designs) but strongly correlated on EGFR (r = 0.64), showing oracle concordance is target-dependent; on EGFR, where 55 of 402 designs have measured binding, gating on concordance lowers binder enrichment rather than raising it.
  \item \textbf{ProtDiff-ESM} \href{https://github.com/brhanufen/protdiff-esm}{(github.com/brhanufen/protdiff-esm)}: Discrete-diffusion protein language model conditioned on frozen ESM-2 with classifier-free guidance toward a fitness target; training and reverse-diffusion sampler built from scratch, evaluated on ProteinGym with a leakage-aware harness whose ESM-2 baseline reproduces the public leaderboard (Pearson r = 0.973).
  \item \textbf{EnhancerDiff} \href{https://github.com/brhanufen/enhancerdiff}{(github.com/brhanufen/enhancerdiff)}: Discrete-diffusion model that designs cell-type-specific 200 bp enhancers (K562, HepG2, SK-N-SH) from the Gosai 2024 MPRA, validated against the Malinois and Enformer oracles.
  \item \textbf{OracleGap} \href{https://github.com/brhanufen/oraclegap}{(github.com/brhanufen/oraclegap)}: Oracle-independence and novelty audit for generative DNA design with Proto (Arc Institute); 93\% of gradient-optimized designs transfer to an independent oracle (agreement $\rho$ = 0.865) with no oracle-hacking detected.
  \item \textbf{SpliceConsensus} \href{https://github.com/brhanufen/spliceconsensus}{(github.com/brhanufen/spliceconsensus)}: Reproducible benchmark of splicing variant-effect predictors against the MFASS experimental assay (27,733 variants); Pangolin ranks best (AUROC 0.888) across five predictors, while 19\% of splice-disrupting variants are missed by every tool, concentrated in the exon interior.
  \item \textbf{AbStab} \href{https://github.com/brhanufen/abstab}{(github.com/brhanufen/abstab)}: Leakage-aware benchmark for antibody thermostability prediction; showed random splits overstate generalization to novel antibodies and that a general PLM (ESM-2) relies on clonal leakage more than an antibody-specific model (AntiBERTy).
  \item \textbf{PerturbVAE} \href{https://github.com/brhanufen/perturbvae}{(github.com/brhanufen/perturbvae)}: Conditional VAE for single-cell perturbation-response prediction, evaluated under random, function-grouped, and combinatorial held-out splits on Norman 2019 and Replogle 2022 Perturb-seq data.
  \item \textbf{NativeReady} \href{https://github.com/brhanufen/nativeready}{(github.com/brhanufen/nativeready)}: Sequence-based triage model and open benchmark for native mass spectrometry suitability (cluster-aware ROC-AUC 0.835 $\pm$ 0.029, n = 635); shipped as a web tool, PyPI SDK, and CLI.
\end{itemize}

% ===== PUBLICATIONS =====
\section{Publications}

\begin{enumerate}[leftmargin=1.5em, topsep=2pt, itemsep=2pt, parsep=0pt]
  \item Ashiquzzaman A*, Lee E*, \textbf{Znabu BF*}, Sakib AN, Chung G, Kim SS, Kim YR, Kwon H-S, Chung E. MoSeq based 3D behavioral profiling uncovers neuropathic behavior changes in diabetic mouse model. \textit{Scientific Reports} 15, 15114 (2025)

  \item \textbf{Znabu BF}, Atif Z. Interpretable deep learning-based multi-omics integration for prognosis in hepatocellular carcinoma. \textit{bioRxiv}. 2026. \href{https://www.biorxiv.org/content/10.64898/2026.04.01.715980v1}{DOI: 10.64898/2026.04.01.715980v1}

  \item \textbf{Znabu BF}, Atif Z. NativeReady: an open benchmark and sequence-based triage model for native mass spectrometry suitability. \textit{bioRxiv}. 2026. \href{https://www.biorxiv.org/content/10.64898/2026.05.03.722506v1}{DOI: 10.64898/2026.05.03.722506v1}

  \item \textbf{Znabu BF}, Atif Z, Devkota P, Alemu RH, K.C. J. A reproducible MFASS benchmark of splice-disruption predictors reveals a shared exon-interior blind spot. \textit{bioRxiv}. 2026. \href{https://www.biorxiv.org/content/10.64898/2026.07.21.739871v1}{DOI: 10.64898/2026.07.21.739871}

  \item \textbf{Znabu BF}, Sexton NR. A genomic foundation model for predicting arbovirus epidemic emergence across RNA virus families. \textit{Bioinformatics} (In preparation)

  \item \textbf{Znabu BF}, Yao Q, Sexton NR. Deciphering the host-range grammar of orthoflaviviruses using foundation model embeddings: a leakage-aware evaluation framework. \textit{PLOS Computational Biology} (In preparation)
\end{enumerate}

% ===== PRESENTATIONS =====
\section{Conference Presentations}

\begin{itemize}[leftmargin=1.2em, topsep=2pt, itemsep=2pt, parsep=0pt]
  \item \textbf{ISMB 2026} (MLCSB COSI), poster accepted: Investigating Host-Range Determinants of Flaviviruses Using Foundation Model Embeddings: A Leakage-Aware Evaluation Framework. Washington, DC (2026)
  \item \textbf{Annual Engineering Symposium}, UNL, poster: Deciphering the host-range grammar of orthoflaviviruses using foundation model embeddings (2026)
  \item \textbf{NCV Annual Virology Symposium}, poster: Cracking the viral code: codon usage and CpG patterns predict host specificity (2025)
\end{itemize}

% ===== TECHNICAL SKILLS =====
\section{Technical Skills}

\textbf{Programming:} Python, R, C++, MATLAB \quad $|$ \quad \textbf{ML/DL:} PyTorch, scikit-learn, XGBoost, Hugging Face Transformers, DNABERT-2, ESM-2, AntiBERTy \quad $|$ \quad \textbf{Infrastructure:} GPU (V100), HPC (Slurm), Git/GitHub \\[3pt]
\textbf{Methods:} Foundation model pre-training (MLM), discrete diffusion, classifier-free guidance, variational autoencoders, LoRA fine-tuning, transfer learning, leakage-aware evaluation, gradient attribution, zero-shot prediction, PCA, UMAP \quad $|$ \quad \textbf{Data:} NCBI GenBank, GISAID, ProteinGym, TCGA, GEO, Biopython, ESMFold, IEDB

% ===== ACADEMIC SERVICE =====
\section{Academic Service}

\entry{Poster Judge}{April 16, 2026}
{School of Biological Sciences Undergraduate Research Symposium, UNL}{Lincoln, NE}

% ===== TEACHING =====
\section{Teaching Experience}

\entry{Lecturer}{Aug 2019 -- Jan 2021}
{Hawassa University}{Ethiopia}
\begin{itemize}[leftmargin=1.2em, topsep=0pt, itemsep=1pt, parsep=0pt]
  \item Taught Biomechanics, Biomaterials, and Biophysics to undergraduate biomedical engineering students.
\end{itemize}

% ===== AWARDS =====
\section{Awards \& Grants}

\begin{itemize}[leftmargin=1.2em, topsep=2pt, itemsep=2pt, parsep=0pt]
  \item Nominated by UNL as one of four students for the Google PhD Fellowship 2026 in AI for Health
  \item Korean Government Scholarship for Master's Study, GIST (2020)
  \item Research Grant for Colostomy Device Development, Hawassa University (2020)
  \item Best B.Sc. Thesis Award, Ethiopian Science, Technology, and Innovation (2018)
\end{itemize}

\end{document}
